\section{MOOSE implementation strategy}
\label{sec:moose_implementation_strategy}

The MOOSE implementation should be built from residual objects that map
directly to the weak forms in Section~\ref{sec:reference_solid_weak_forms}.  The
first implementation milestone is not broad feature coverage; it is one
validated residual path on the solid reference configuration.

\subsection{Proposed object responsibilities}
\label{sec:object_responsibilities}

\begin{table}[h]
	\centering
	\caption{Initial implementation interfaces.  Names are descriptive
	placeholders, not required C++ class names.}
	\label{tab:implementation_interfaces}
	\small
	\begin{tabular}{p{0.23\linewidth}p{0.34\linewidth}p{0.26\linewidth}}
		\toprule
		Object role & Computes & Weak-form target \\
		\midrule
		Solid kinematics material & \(\mbf F\), \(J\), \(\mbf F^{-1}\), \(\mbf v_s\) & all solid-reference residuals \\
		Phase advection material & \(\mbf v_f\) and \(D_f(\cdot)/Dt\) from \(\mbf W_f\) & Eq.~\eqref{eq:fluid_material_derivative_on_solid_frame} \\
		Thermodynamics material & \(\bar p_f\), \(\mu_\xi^\alpha\), \(\bar p_E\), \(B\), \(\pi_\xi^\alpha\), \(\Pi_\xi\), \(\lambda\) & flux, momentum, algebraic constraints \\
		Reaction material & \(\dot r_{(m)}\), \(\dot c_\xi^\alpha\), affinities & source and conversion terms \\
		Reference flux material & \(\mbf W_f\), \(\widetilde{\mbf K}_f\) & Eq.~\eqref{eq:fluid_component_residual} \\
		Stress material & \(\mbf P^{\prime\prime}\), pressure stress split & Eq.~\eqref{eq:momentum_residual} \\
		Component kernels & storage, flux, source residuals & Eq.~\eqref{eq:fluid_component_residual} \\
		Solid kernels & solid storage and conversion residuals & Eq.~\eqref{eq:solid_component_residual} \\
		Conversion-potential kernels & solid-frame evolution of \(\mu_\xi^\alpha\) & Eqs.~\eqref{eq:fluid_mu_residual}--\eqref{eq:solid_mu_residual} \\
		Constraint kernels/material solves & partition, pressure, and composition relations & Eqs.~\eqref{eq:implementation_partition_relations}--\eqref{eq:implementation_pi_difference_system} \\
		Momentum kernels & skeleton force balance & Eq.~\eqref{eq:momentum_residual} \\
		\bottomrule
	\end{tabular}
\end{table}

All objects should use automatic differentiation from the start.  Non-AD helper
properties may be acceptable only for constants or diagnostics that do not enter
Newton derivatives.

\subsection{Automatic-differentiation boundary}
\label{sec:ad_boundary}

The implementation should separate scalar thermodynamic differentiation from
tensorial kinematics and residual assembly.  The current MOOSE finding is that
\texttt{ADDerivativeParsedMaterial} can represent scalar Helmholtz-free-energy
closures and expose derivatives such as pressures, chemical potentials, and
selected second derivatives when the independent scalar arguments are chosen
explicitly.  That capability is useful for local thermodynamic derivatives, but
it does not remove the need for explicit AD C++ code for solid and fluid
deformation-gradient kinematics, Piola transformations, pulled-back modified
permeability, source closures, constraint residuals, conversion-potential
evolution, and the weak forms in
Section~\ref{sec:reference_solid_weak_forms}.

Consequently, parsed derivative materials should be treated as thermodynamic
building blocks rather than as the implementation architecture.  Kernels and
materials that compute \(\mbf F\), \(J\), \(\mbf F^{-1}\), \(\mbf W_f\),
\(\mbf v_f\), \(\mbf P^{\prime\prime}\),
\(B\bar p_EJ\mbf F^{-T}\), and the storage/source terms must retain full AD
dependence on the finite-element variables and their gradients.

This boundary gives a practical rule for the first code review: a parsed free
energy may define a scalar potential and its partial derivatives, but a MOOSE
kernel must still show exactly which weak-form term it contributes to.  If a
derivative is produced by a parsed material, the input file must document the
independent variables and the maximum derivative order requested.  If a
derivative is produced by compiled C++, the material header or source should
state the thermodynamic identity it implements.

\subsection{Primitive variables}
\label{sec:primitive_variables}

The scaffold leaves the primitive variables open.  The choice should be made by
checking which variables give a nonsingular thermodynamic Jacobian and a clean
mapping to standard reservoir limits.  Candidate sets include:
\begin{itemize}[leftmargin=*]
	\item displacement plus independent Lagrangian component masses;
	\item displacement, phase pressures, saturations, and component fractions with algebraic closures;
	\item displacement plus thermodynamic potentials and constraint multipliers.
\end{itemize}

Whichever set is chosen, materials must provide the physical quantities used in
the residuals: \(M_f^\alpha\), \(M_s\), \(\mbf W_f\),
\(\mbf P^{\prime\prime}\), \(\bar p_E\), \(B\), \(\mu_\xi^\alpha\),
\(\pi_\xi^\alpha\), \(\Pi_\xi\), \(\lambda\), and \(\dot r_{(m)}\).

The first reviewable implementation plan is closed only for the restricted
problem stated in Section~\ref{sec:closure_requirement}.  The broader
primitive-variable contract remains open; once chosen, it must be recorded as
the map from finite-element variables to the physical quantities appearing in
the residuals.

\subsection{Discretization defaults}
\label{sec:discretization_defaults}

The first finite-element discretization should use continuous Lagrange elements
with a mixed order sufficient to avoid locking or pressure oscillations in the
poroelastic limit.  A conservative starting point is quadratic displacement and
linear scalar variables, matching the Taylor--Hood pattern already used in the
three-phase implementation.  Stabilization should be introduced only after a
specific failure mode is identified and tied to a documented weak-form
approximation.

\subsection{Traceability rules}
\label{sec:traceability_rules}

Each eventual MOOSE object should cite the controlling equation in this paper
and the corresponding equation or section in the theory manuscript.  Each input
deck should state the reduction it exercises: compositional flow, black-oil
limit, poromechanics, reaction-only conversion, phase equilibrium, or a coupled
benchmark.  Generated exodus output is not validation by itself; every
verification case needs an observable, a reference result, and a tolerance.
